\section*{Which existing result could explain a win?}
\textbf{Learned read, then read again.} End-to-End Memory Networks (Section 2) combine a weighted read with the query and use it for the next hop. A single register is sufficient for our tables. Recurrent Independent Mechanisms already maintain multiple private states and learn when to communicate. The algebra on page 2 places the present fresh arm inside the model class of a tied block-recurrent network. A wiring diagram alone provides no novelty. \cite{mem,rims}
\par
\textbf{Reweight, select or synthesize.} Multiplex Thinking samples several token embeddings, merges them and takes another LM step. GRAM (Sections 2.1--2.3) draws recursive stochastic latents from an inference prior and can apply a scorer or voting across trajectories. Guided Reasoning (Sections 4.1--4.3) resamples particles online using a learned stopping head. Both offline selection and online allocation are relevant controls. On ARC, Latent Program Networks (Sections 4.1--4.3) fit a shared latent program to observed demonstrations by sampling or optimization; charge their program evaluations and gradients. \cite{multiplex,gram,guided,lpn}
\par
\textbf{Native capability comparisons.} TRM's released implementation carries answer and workspace states; its stock evaluation forces the configured maximum depth. The ARC evaluator also uses puzzle identifiers and augmented predictions. FPRM means Fixed-Point Reasoning Model; its release uses the name FPTRM and a damped fixed-point update with convergence-based stopping. InfiLoop aggregates history within a single trajectory. Its public Sudoku solver retries surviving boards from scratch at different depths and checks a legal puzzle certificate. An equal-depth run and a native solver run answer different questions. SE-RRM supplies symbol equivariance. Its Sudoku Table 2 varies the number of supervision steps; Section 3.2 describes its ARC demonstration exposure. \cite{trm,fprm,infi,se}
\par
These releases have different task interfaces, conditioning, stopping and selection costs. None is a four-symbol PH1 or PH2 checkpoint. The ARC bridge used a fixed untrained program on the public training set; comparing its 2.75\% to an ARC evaluation-set headline would be invalid. We ran no released checkpoint in this revision. Before any natural benchmark claim, train or run the strongest compatible models at their native settings and with common legal information, report their costs, and include a trained single-trajectory rival. The accompanying \texttt{prior-art.md} records source sections, checkpoint paths, reviewer discussion and access limitations.
\par
\textbf{Resource decision.} Retire these two performance claims at the tested budgets. A repeat of the conditional diagnostic on one GH200 is technically prepared (15 minute wall limit, at most 0.25 requested GPU-hours); it would measure GPU dispatch and numerical parity only. The recommended allocation is zero. A future study needs an open answer space and a genuine future operation that can be causally changed by complementary consequences. No cluster job or spending commitment is pending.
